\documentclass[12pt]{article}
\usepackage[margin=1in]{geometry}
\usepackage{graphicx}
\usepackage[hidelinks]{hyperref}
\usepackage{float}

\title{Global Electricity Access and Renewable Energy: Key Patterns}
\author{Grace Carver}
\date{May 2026}

\pagestyle{empty}

\begin{document}

\vspace*{-2cm}
\maketitle
\vspace{-1.5cm}

\section*{Overview}

Electricity access and renewable energy adoption are foundational to economic development and climate action. This brief summarizes global patterns in electricity access and the adoption of renewable energy sources across countries, with emphasis on disparities by income level and region. Data are drawn from the World Bank's World Development Indicators (2021 cross-section; 2000–2021 time series).

\section*{Income and Electricity Access}

\begin{figure}[H]
\centering
\includegraphics[width=0.9\linewidth]{2026_Application_Day_3/output/figures/day4/fig1_access_vs_gdp_scatter.png}
\end{figure}

Electricity access exhibits a strong income gradient. Countries with GDP per capita below \$5,000 show wide variation in access (10\% to 100\%), while nearly all countries above \$20,000 per capita achieve near-universal access. Sub-Saharan Africa lags significantly: most countries in this region cluster in the lower-left (low income, incomplete access), whereas Europe and North America occupy the upper-right (high income, high access). This pattern reflects both the capital requirements for electrification and the concentration of economic activity in electrified areas. Notably, South and East Asia shows rapid progress: several countries have escaped the income-access trap through targeted infrastructure investment.

\section*{Income and Renewable Energy}

\begin{figure}[H]
\centering
\includegraphics[width=0.9\linewidth]{2026_Application_Day_3/output/figures/day4/fig2_renewable_vs_gdp_scatter.png}
\end{figure}

Renewable electricity adoption shows less clear dependence on income than access does. While some high-income countries cluster at high renewable shares (particularly Europe and North America, driven by hydropower, wind, and solar), others remain below 20\% renewable despite wealth. Conversely, some low-income countries achieve substantial renewable shares, typically through abundant hydropower (e.g., Sub-Saharan Africa) or biomass. This suggests that renewable adoption is driven more by geography and resource endowment than by income alone. Regional variation is pronounced: Latin America and MENA show significant renewable capacity, often from hydropower and solar resources.

\section*{Regional Trends: Pace of Expansion}

\begin{figure}[H]
\centering
\includegraphics[width=0.9\linewidth]{2026_Application_Day_3/output/figures/day4/fig8_access_regional_facets.png}
\end{figure}

Electricity access has expanded globally since 2000, but the pace and current level vary markedly by region. Europe and North America began above 95\% access and have reached saturation. Latin America and MENA saw modest gains from already-high baselines (85\% to 95\%), while South and East Asia achieved the steepest growth: rising from roughly 75\% in 2000 to near 95\% in 2021. Sub-Saharan Africa, starting from the lowest baseline (around 25\%), has expanded access to approximately 50\% over the two decades—substantial progress in absolute terms, but the slowest pace of improvement and the largest remaining access gap. This regional divergence reflects both historical infrastructure investment and current capacity constraints.

\section*{Sources}

World Bank. (2023). \textit{World Development Indicators}. Retrieved from \url{https://data.worldbank.org/}. Data accessed May 2026. License: CC BY 4.0.

\smallskip

Indicators: Electricity access (share of population, \%); Renewable electricity output (share of total electricity output, \%); GDP per capita (current USD, log scale).

\end{document}